\documentclass[10pt,a4paper]{amsart}
\usepackage{amsmath,amssymb,amsthm,graphicx}
\usepackage{verbatim}
\newenvironment{reference}{\comment}{\endcomment}
%\newenvironment{reference}{}{}
\newenvironment{slogan}{\comment}{\endcomment}
\newenvironment{history}{\comment}{\endcomment}

% For commutative diagrams we use Xy-pic
\usepackage[all]{xy}

% We use 2cell for 2-commutative diagrams.
\xyoption{2cell}
% Theorem environments.
%
\theoremstyle{plain}
\newtheorem{theorem}[subsection]{Theorem}
\newtheorem{proposition}[subsection]{Proposition}
\newtheorem{lemma}[subsection]{Lemma}

\theoremstyle{definition}
\newtheorem{definition}[subsection]{Definition}
\newtheorem{example}[subsection]{Example}
\newtheorem{exercise}[subsection]{Exercise}
\newtheorem{situation}[subsection]{Situation}

\theoremstyle{remark}
\newtheorem{remark}[subsection]{Remark}
\newtheorem{remarks}[subsection]{Remarks}

\numberwithin{equation}{subsection}

% Macros
%
\def\lim{\mathop{\mathrm{lim}}\nolimits}
\def\colim{\mathop{\mathrm{colim}}\nolimits}
\def\Spec{\mathop{\mathrm{Spec}}}
\def\Hom{\mathop{\mathrm{Hom}}\nolimits}
\def\Ext{\mathop{\mathrm{Ext}}\nolimits}
\def\SheafHom{\mathop{\mathcal{H}\!\mathit{om}}\nolimits}
\def\SheafExt{\mathop{\mathcal{E}\!\mathit{xt}}\nolimits}
\def\Sch{\mathit{Sch}}
\def\Mor{\mathop{\mathrm{Mor}}\nolimits}
\def\Ob{\mathop{\mathrm{Ob}}\nolimits}
\def\Sh{\mathop{\mathit{Sh}}\nolimits}
\def\NL{\mathop{N\!L}\nolimits}
\def\CH{\mathop{\mathrm{CH}}\nolimits}
\def\proetale{{pro\text{-}\acute{e}tale}}
\def\etale{{\acute{e}tale}}
\def\QCoh{\mathit{QCoh}}
\def\Ker{\mathop{\mathrm{Ker}}}
\def\Im{\mathop{\mathrm{Im}}}
\def\Coker{\mathop{\mathrm{Coker}}}
\def\Coim{\mathop{\mathrm{Coim}}}

% Boxtimes
%
\DeclareMathSymbol{\boxtimes}{\mathbin}{AMSa}{"02}

%
% Macros for moduli stacks/spaces
%
\def\QCohstack{\mathcal{QC}\!\mathit{oh}}
\def\Cohstack{\mathcal{C}\!\mathit{oh}}
\def\Spacesstack{\mathcal{S}\!\mathit{paces}}
\def\Quotfunctor{\mathrm{Quot}}
\def\Hilbfunctor{\mathrm{Hilb}}
\def\Curvesstack{\mathcal{C}\!\mathit{urves}}
\def\Polarizedstack{\mathcal{P}\!\mathit{olarized}}
\def\Complexesstack{\mathcal{C}\!\mathit{omplexes}}
% \Pic is the operator that assigns to X its picard group, usage \Pic(X)
% \Picardstack_{X/B} denotes the Picard stack of X over B
% \Picardfunctor_{X/B} denotes the Picard functor of X over B
\def\Pic{\mathop{\mathrm{Pic}}\nolimits}
\def\Picardstack{\mathcal{P}\!\mathit{ic}}
\def\Picardfunctor{\mathrm{Pic}}
\def\Deformationcategory{\mathcal{D}\!\mathit{ef}}
\usepackage{cleveref}
\crefname{originalreference}{original reference}{original references}
\Crefname{originalreference}{Original reference}{Original references}
\makeatletter
\newlabel{spaces-divisors-lemma-ample-base-change}{{0D33}{source}{}{sourcepointer-1}{}}
\newlabel{spaces-divisors-lemma-ample-base-change@cref}{{[originalreference][0][]0D33}{[1][0][]source}{}{}{}}
\makeatother
\crefname{originalreference}{original reference}{original references}
\Crefname{originalreference}{Original reference}{Original references}
\makeatletter
\newlabel{spaces-divisors-lemma-relatively-ample-local}{{0D36}{source}{}{sourcepointer-2}{}}
\newlabel{spaces-divisors-lemma-relatively-ample-local@cref}{{[originalreference][0][]0D36}{[1][0][]source}{}{}{}}
\makeatother
\crefname{originalreference}{original reference}{original references}
\Crefname{originalreference}{Original reference}{Original references}
\makeatletter
\newlabel{spaces-topologies-lemma-refine-fpqc-schemes}{{0419}{source}{}{sourcepointer-3}{}}
\newlabel{spaces-topologies-lemma-refine-fpqc-schemes@cref}{{[originalreference][0][]0419}{[1][0][]source}{}{}{}}
\makeatother
\crefname{originalreference}{original reference}{original references}
\Crefname{originalreference}{Original reference}{Original references}
\makeatletter
\newlabel{spaces-divisors-lemma-relatively-ample-properties}{{0D34}{source}{}{sourcepointer-4}{}}
\newlabel{spaces-divisors-lemma-relatively-ample-properties@cref}{{[originalreference][0][]0D34}{[1][0][]source}{}{}{}}
\makeatother
\crefname{originallemma}{original lemma}{original lemmas}
\Crefname{originallemma}{Original lemma}{Original lemmas}
\makeatletter
\newlabel{lemma-descending-property-quasi-compact}{{041L}{source}{}{sourcepointer-5}{}}
\newlabel{lemma-descending-property-quasi-compact@cref}{{[originallemma][0][]041L}{[1][0][]source}{}{}{}}
\makeatother
\crefname{originallemma}{original lemma}{original lemmas}
\Crefname{originallemma}{Original lemma}{Original lemmas}
\makeatletter
\newlabel{lemma-descending-property-separated}{{0421}{source}{}{sourcepointer-6}{}}
\newlabel{lemma-descending-property-separated@cref}{{[originallemma][0][]0421}{[1][0][]source}{}{}{}}
\makeatother
\crefname{originalreference}{original reference}{original references}
\Crefname{originalreference}{Original reference}{Original references}
\makeatletter
\newlabel{spaces-morphisms-lemma-pushforward}{{03M9}{source}{}{sourcepointer-7}{}}
\newlabel{spaces-morphisms-lemma-pushforward@cref}{{[originalreference][0][]03M9}{[1][0][]source}{}{}{}}
\makeatother
\crefname{originalreference}{original reference}{original references}
\Crefname{originalreference}{Original reference}{Original references}
\makeatletter
\newlabel{spaces-cohomology-lemma-flat-base-change-cohomology}{{073K}{source}{}{sourcepointer-8}{}}
\newlabel{spaces-cohomology-lemma-flat-base-change-cohomology@cref}{{[originalreference][0][]073K}{[1][0][]source}{}{}{}}
\makeatother
\crefname{originalreference}{original reference}{original references}
\Crefname{originalreference}{Original reference}{Original references}
\makeatletter
\newlabel{morphisms-lemma-characterize-relatively-ample}{{01VJ}{source}{}{sourcepointer-9}{}}
\newlabel{morphisms-lemma-characterize-relatively-ample@cref}{{[originalreference][0][]01VJ}{[1][0][]source}{}{}{}}
\makeatother
\crefname{originalreference}{original reference}{original references}
\Crefname{originalreference}{Original reference}{Original references}
\makeatletter
\newlabel{spaces-divisors-lemma-invertible-map-into-relative-proj}{{0D2Z}{source}{}{sourcepointer-10}{}}
\newlabel{spaces-divisors-lemma-invertible-map-into-relative-proj@cref}{{[originalreference][0][]0D2Z}{[1][0][]source}{}{}{}}
\makeatother
\crefname{originallemma}{original lemma}{original lemmas}
\Crefname{originallemma}{Original lemma}{Original lemmas}
\makeatletter
\newlabel{lemma-descending-property-open-immersion}{{041X}{source}{}{sourcepointer-11}{}}
\newlabel{lemma-descending-property-open-immersion@cref}{{[originallemma][0][]041X}{[1][0][]source}{}{}{}}
\makeatother
\title{Ampleness is fpqc local on the base}
\author{The Stacks Project Authors}
\date{stacks-project:spaces-descent.tex:local-snapshot-20261012; GNU FDL 1.2 or later}
\begin{document}
\maketitle
\noindent\textit{Editorial note.} All delivered mathematical text is the author's own, copied verbatim from the pinned Stacks Project chapter file spaces-descent.tex; no statement or proof step was rewritten, reordered or removed.\par\medskip
\noindent\textit{Editorial note.} The counted span cites items from Divisors on Spaces, Topologies on Spaces, Morphisms of Spaces, Cohomology of Spaces and Morphisms, plus same-chapter descent lemmas (quasi-compact, separated, open immersion), recorded in assembly.external\_references with canonical tags as navigation pointers.\par\medskip
\noindent\textit{Allowed standard prerequisites (editorial).}\begin{itemize}
\item Base change of relatively ample modules. Stacks Project, Divisors on Spaces, Lemma spaces-divisors-lemma-ample-base-change. Proves (1) implies (2).
\item Relative ampleness is \textbackslash{}'etale local on the base. Stacks Project, Divisors on Spaces, Lemma spaces-divisors-lemma-relatively-ample-local. Lets the proof assume \$Y\$ is a scheme.
\item Refining fpqc coverings by schemes. Stacks Project, Topologies on Spaces, Lemma spaces-topologies-lemma-refine-fpqc-schemes. Lets the proof assume each \$Y\_i\$ is a scheme.
\item Properties of relatively ample morphisms. Stacks Project, Divisors on Spaces, Lemma spaces-divisors-lemma-relatively-ample-properties. Gives representability, quasi-compactness and separatedness of \$X\_i \textbackslash{}to Y\_i\$, and the final ampleness characterization.
\item Pushforward of quasi-coherent sheaves. Stacks Project, Morphisms of Spaces, Lemma spaces-morphisms-lemma-pushforward. Shows \$\textbackslash{}mathcal\{A\}\$ is a quasi-coherent graded \$\textbackslash{}mathcal\{O\}\_Y\$-algebra.
\item Flat base change of pushforwards. Stacks Project, Cohomology of Spaces, Lemma spaces-cohomology-lemma-flat-base-change-cohomology. Shows the formation of \$\textbackslash{}mathcal\{A\}\$ commutes with flat base change, so \$\textbackslash{}mathcal\{A\}\_i = g\_i\textasciicircum{}*\textbackslash{}mathcal\{A\}\$.
\item Map into relative Proj from an invertible module. Stacks Project, Divisors on Spaces, Lemma spaces-divisors-lemma-invertible-map-into-relative-proj. Supplies the open subset \$U(\textbackslash{}psi)\$ and its behavior under base change.
\item Characterization of relatively ample invertible modules. Stacks Project, Morphisms, Lemma morphisms-lemma-characterize-relatively-ample. Supplies the stalk surjectivity criterion and the Proj open-immersion criterion.
\item fpqc descent of open immersions. Stacks Project, Descent on Spaces, Lemma lemma-descending-property-open-immersion. Descends the open immersion \$r\_\{\textbackslash{}mathcal\{L\}, \textbackslash{}psi\}\$ from the \$Y\_i\$.
\end{itemize}\medskip
\ifdefined\hypertarget\hypertarget{sourcepointer-1}{}\fi
\noindent\textit{Source reference.} Original reference 0D33: sources/stacks-chow/tags.txt, source lines 16403--16403. Cross-chapter standard prerequisite cited by the author inside the counted span; it is used as a pointer and its own proof is not part of this item.\par\medskip
\ifdefined\hypertarget\hypertarget{sourcepointer-2}{}\fi
\noindent\textit{Source reference.} Original reference 0D36: sources/stacks-chow/tags.txt, source lines 16406--16406. Cross-chapter standard prerequisite cited by the author inside the counted span; it is used as a pointer and its own proof is not part of this item.\par\medskip
\ifdefined\hypertarget\hypertarget{sourcepointer-3}{}\fi
\noindent\textit{Source reference.} Original reference 0419: sources/stacks-chow/tags.txt, source lines 5284--5284. Cross-chapter standard prerequisite cited by the author inside the counted span; it is used as a pointer and its own proof is not part of this item.\par\medskip
\ifdefined\hypertarget\hypertarget{sourcepointer-4}{}\fi
\noindent\textit{Source reference.} Original reference 0D34: sources/stacks-chow/tags.txt, source lines 16404--16404. Cross-chapter standard prerequisite cited by the author inside the counted span; it is used as a pointer and its own proof is not part of this item.\par\medskip
\ifdefined\hypertarget\hypertarget{sourcepointer-5}{}\fi
\noindent\textit{Source reference.} Original lemma 041L: sources/stacks-chow/tags.txt, source lines 5296--5296. Cross-chapter standard prerequisite cited by the author inside the counted span; it is used as a pointer and its own proof is not part of this item.\par\medskip
\ifdefined\hypertarget\hypertarget{sourcepointer-6}{}\fi
\noindent\textit{Source reference.} Original lemma 0421: sources/stacks-chow/tags.txt, source lines 5312--5312. Cross-chapter standard prerequisite cited by the author inside the counted span; it is used as a pointer and its own proof is not part of this item.\par\medskip
\ifdefined\hypertarget\hypertarget{sourcepointer-7}{}\fi
\noindent\textit{Source reference.} Original reference 03M9: sources/stacks-chow/tags.txt, source lines 4741--4741. Cross-chapter standard prerequisite cited by the author inside the counted span; it is used as a pointer and its own proof is not part of this item.\par\medskip
\ifdefined\hypertarget\hypertarget{sourcepointer-8}{}\fi
\noindent\textit{Source reference.} Original reference 073K: sources/stacks-chow/tags.txt, source lines 9064--9064. Cross-chapter standard prerequisite cited by the author inside the counted span; it is used as a pointer and its own proof is not part of this item.\par\medskip
\ifdefined\hypertarget\hypertarget{sourcepointer-9}{}\fi
\noindent\textit{Source reference.} Original reference 01VJ: sources/stacks-chow/tags.txt, source lines 2448--2448. Cross-chapter standard prerequisite cited by the author inside the counted span; it is used as a pointer and its own proof is not part of this item.\par\medskip
\ifdefined\hypertarget\hypertarget{sourcepointer-10}{}\fi
\noindent\textit{Source reference.} Original reference 0D2Z: sources/stacks-chow/tags.txt, source lines 16399--16399. Cross-chapter standard prerequisite cited by the author inside the counted span; it is used as a pointer and its own proof is not part of this item.\par\medskip
\ifdefined\hypertarget\hypertarget{sourcepointer-11}{}\fi
\noindent\textit{Source reference.} Original lemma 041X: sources/stacks-chow/tags.txt, source lines 5308--5308. Cross-chapter standard prerequisite cited by the author inside the counted span; it is used as a pointer and its own proof is not part of this item.\par\medskip
\begin{lemma}
\label{lemma-descending-property-ample}
Let $S$ be a scheme.
Let $f : X \to Y$ be a morphism of algebraic spaces over $S$.
Let $\mathcal{L}$ be an invertible $\mathcal{O}_X$-module.
Let $\{g_i : Y_i \to Y\}_{i \in I}$ be an fpqc covering.
Let $f_i : X_i \to Y_i$ be the base change of $f$ and let $\mathcal{L}_i$
be the pullback of $\mathcal{L}$ to $X_i$.
The following are equivalent
\begin{enumerate}
\item $\mathcal{L}$ is ample on $X/Y$, and
\item $\mathcal{L}_i$ is ample on $X_i/Y_i$
for every $i \in I$.
\end{enumerate}
\end{lemma}

\begin{proof}
The implication (1) $\Rightarrow$ (2) follows from
Divisors on Spaces, Lemma \ref{spaces-divisors-lemma-ample-base-change}.
Assume (2). To check $\mathcal{L}$ is ample on $X/Y$ we may
work \'etale locally on $Y$, see
Divisors on Spaces, Lemma \ref{spaces-divisors-lemma-relatively-ample-local}.
Thus we may assume that $Y$ is a scheme and then we may
in turn assume each $Y_i$ is a scheme too, see
Topologies on Spaces, Lemma \ref{spaces-topologies-lemma-refine-fpqc-schemes}.
In other words, we may assume that
$\{Y_i \to Y\}$ is an fpqc covering of schemes.

\medskip\noindent
By Divisors on Spaces, Lemma
\ref{spaces-divisors-lemma-relatively-ample-properties}
we see that $X_i \to Y_i$ is representable (i.e., $X_i$ is a scheme),
quasi-compact, and separated. Hence $f$ is quasi-compact and separated by
Lemmas \ref{lemma-descending-property-quasi-compact} and
\ref{lemma-descending-property-separated}.
This means that
$\mathcal{A} = \bigoplus_{d \geq 0} f_*\mathcal{L}^{\otimes d}$
is a quasi-coherent graded $\mathcal{O}_Y$-algebra
(Morphisms of Spaces, Lemma \ref{spaces-morphisms-lemma-pushforward}).
Moreover, the formation of $\mathcal{A}$ commutes with flat
base change by
Cohomology of Spaces, Lemma
\ref{spaces-cohomology-lemma-flat-base-change-cohomology}.
In particular, if we set
$\mathcal{A}_i = \bigoplus_{d \geq 0} f_{i, *}\mathcal{L}_i^{\otimes d}$
then we have $\mathcal{A}_i = g_i^*\mathcal{A}$.
It follows that the natural maps
$\psi_d : f^*\mathcal{A}_d \to \mathcal{L}^{\otimes d}$
of $\mathcal{O}_X$
pullback to give the natural maps
$\psi_{i, d} : f_i^*(\mathcal{A}_i)_d \to \mathcal{L}_i^{\otimes d}$
of $\mathcal{O}_{X_i}$-modules. Since $\mathcal{L}_i$ is ample on $X_i/Y_i$
we see that for any point $x_i \in X_i$, there exists a $d \geq 1$
such that $f_i^*(\mathcal{A}_i)_d \to \mathcal{L}_i^{\otimes d}$
is surjective on stalks at $x_i$. This follows either directly
from the definition of a relatively ample module or from
Morphisms, Lemma \ref{morphisms-lemma-characterize-relatively-ample}.
If $x \in |X|$, then we can choose an $i$ and an $x_i \in X_i$
mapping to $x$. Since
$\mathcal{O}_{X, \overline{x}} \to \mathcal{O}_{X_i, \overline{x}_i}$
is flat hence faithfully flat, we conclude that for every $x \in |X|$
there exists a $d \geq 1$ such that
$f^*\mathcal{A}_d \to \mathcal{L}^{\otimes d}$
is surjective on stalks at $x$.
This implies that the open subset $U(\psi) \subset X$ of
Divisors on Spaces, Lemma
\ref{spaces-divisors-lemma-invertible-map-into-relative-proj}
corresponding to the map
$\psi : f^*\mathcal{A} \to \bigoplus_{d \geq 0} \mathcal{L}^{\otimes d}$
of graded $\mathcal{O}_X$-algebras
is equal to $X$. Consider the corresponding morphism
$$
r_{\mathcal{L}, \psi} : X \longrightarrow \underline{\text{Proj}}_Y(\mathcal{A})
$$
It is clear from the above that the base change of
$r_{\mathcal{L}, \psi}$ to $Y_i$ is the morphism
$r_{\mathcal{L}_i, \psi_i}$ which is an open immersion by
Morphisms, Lemma \ref{morphisms-lemma-characterize-relatively-ample}.
Hence $r_{\mathcal{L}, \psi}$ is an open immersion
by Lemma \ref{lemma-descending-property-open-immersion}.
Hence $X$ is a scheme
and we conclude $\mathcal{L}$ is ample on $X/Y$ by
Morphisms, Lemma \ref{morphisms-lemma-characterize-relatively-ample}.
\end{proof}

\end{document}
